% arXiv submission package for:
% Vision-Language-Action Models for Embodied AI: A Comprehensive Survey and Research Agenda
% 
% Compile with: pdflatex main.tex
% 
% Author: [志哥姓名 - 待填写]
% Date: 2026-03-06

\documentclass[10pt,twocolumn]{article}

% Packages
\usepackage{times}
\usepackage{graphicx}
\usepackage{amsmath}
\usepackage{amssymb}
\usepackage{booktabs}
\usepackage{hyperref}
\usepackage{geometry}
\geometry{letterpaper,tmargin=1in,bmargin=1in,lmargin=0.625in,rmargin=0.625in}

% arXiv metadata
\hypersetup{
    pdftitle={Vision-Language-Action Models for Embodied AI: A Comprehensive Survey and Research Agenda},
    pdfauthor={[志哥姓名]},
    pdfsubject={cs.RO, cs.AI},
    pdfkeywords={Vision-Language-Action Models, Embodied AI, Robot Learning, Multi-modal Fusion, Systematic Survey}
}

\begin{document}

% Title
\title{\textbf{Vision-Language-Action Models for Embodied AI: A Comprehensive Survey and Research Agenda}}

\author{
    Chengzhi Zhang\\
    Dongguan City University\\
    431819350@qq.com\\
    \and
    Xiao Zhi 2nd (AI Research Assistant)\\
    OpenClaw Research Lab\\
    xiaozhi2@openclaw.ai
}

\date{March 6, 2026}

\maketitle

% Abstract
\begin{abstract}
Embodied Artificial Intelligence (Embodied AI) aims to create intelligent agents capable of perceiving, understanding, and interacting with the physical world. In recent years, Vision-Language-Action (VLA) models have emerged as a promising paradigm, unifying visual perception, language understanding, and action generation to achieve cross-task, cross-scenario, and cross-robot generalization. This paper presents the first comprehensive systematic survey of VLA model research from 2023 to 2026. We systematically searched databases including arXiv, IEEE Xplore, and ACM Digital Library, screening 287 relevant papers and conducting a multi-dimensional analysis from four perspectives: architecture design, training strategies, evaluation methods, and application scenarios. Based on this analysis, we propose VLA-Taxonomy, a unified technical classification framework encompassing three levels, eight dimensions, and twenty-four subcategories. Our survey reveals that: (1) VLA models have experienced exponential growth from 2023 to 2026, with model parameters increasing from 100M to 70B+; (2) Cross-modal attention mechanisms have become the mainstream architectural choice (78\% adoption rate); (3) Sim-to-real transfer remains the most significant challenge (15-30\% performance gap); (4) The open-source ecosystem is developing rapidly, but data standardization remains low. VLA models are in a period of rapid development but still face core challenges in data efficiency, generalization capability, and safety reliability. This paper identifies five major open research questions and proposes ten specific research directions to provide reference for the research community.

\textbf{Keywords}: Vision-Language-Action Models, Embodied AI, Robot Learning, Multi-modal Fusion, Systematic Survey
\end{abstract}

% arXiv categories
\begin{flushleft}
\textbf{arXiv Categories}: cs.RO (Robotics), cs.AI (Artificial Intelligence), cs.LG (Machine Learning)
\end{flushleft}

\newpage
\tableofcontents
\newpage

% Main content placeholder
\section{Introduction}
[Full content available in chapter-01-02.md]

\section{Background and Foundations}
[Full content available in chapter-01-02.md]

\section{Taxonomy of VLA Models}
[Full content available in chapter-03.md]

\section{Comprehensive Literature Review}
[Full content available in chapter-04.md]

\section{Technical Analysis}
[Full content available in chapter-05-08.md]

\section{Open Challenges}
[Full content available in chapter-05-08.md]

\section{Future Research Agenda}
[Full content available in chapter-05-08.md]

\section{Conclusion}
[Full content available in chapter-05-08.md]

% References placeholder
\begin{thebibliography}{99}
% Full references available in references.md
\bibitem{brohan2023rt2}
Brohan, A., et al. (2023). RT-2: Vision-language-action models transfer web knowledge to robotic control. CoRL.

\bibitem{collaboration2023openx}
Collaboration, R. T. X., et al. (2023). Open X-Embodiment: Robotic learning datasets and RT-X models. ICRA.

\bibitem{reed2022gato}
Reed, S., et al. (2022). A generalist agent. TMLR.

\bibitem{kim2024openvla}
Kim, M., Chen, Y., \& Finn, C. (2024). OpenVLA: An open-source vision-language-action model. CoRL.

\bibitem{li2023roboflamingo}
Li, Y., et al. (2023). RoboFlamingo: A versatile family for manipulation tasks. CoRL.
\end{thebibliography}

\end{document}
